\documentclass[reqno,11pt]{amsart}
\usepackage{math327-lecture}
\usepackage{needspace}
\usepackage{tikz}
\lectureinfo{21}{Oct. 28, 2026}{Group Actions, Orbits, Stabilizers, and Rotations}

% Draft based on handwritten Lectures 15 (Oct. 6, 2025) and
% 16 (Oct. 8, 2025). Retain the full substantive sequence for editing.
\begin{document}
\makelecturetitle

\section{Group actions}

Groups become especially useful when they act on other objects.
Let $G$ be a group and $S$ a set. When studying the correspondence theorem, we introduce the following notion:

\begin{define}[Left group action]\label{def:action}
A \emph{left action} of $G$ on $S$ is a map
\[
  G\times S\longrightarrow S,\qquad (g,s)\longmapsto g\cdot s,
\]
such that, for all $g,h\in G$ and $s\in S$,
\[
  1\cdot s=s,
  \qquad (gh)\cdot s=g\cdot(h\cdot s).
\]
We also write $G\curvearrowright S$.
\end{define}

\begin{example}[Permutations]
The symmetric group $S_n$ acts on $S=\{1,\ldots,n\}$ by
$\sigma\cdot j=\sigma(j)$. The identity fixes each $j$, and the
action axiom is the usual composition rule
$(\sigma\tau)(j)=\sigma(\tau(j))$.
\end{example}

\begin{example}[Invertible matrices]
The group $\GL_n(\R)$ acts on $\R^n$ by matrix multiplication:
$A\cdot v=Av$. The identity matrix fixes every vector, and
$(AB)v=A(Bv)$.
\end{example}

\begin{remark}
Each $g\in G$ induces a bijection
$m_g:S\to S$, $s\mapsto g\cdot s$. Its inverse is $m_{g^{-1}}$,
since
\[
  g^{-1}\cdot(g\cdot s)=(g^{-1}g)\cdot s=s.
\]
The action law says $m_{gh}=m_g\circ m_h$. Thus an action gives a
homomorphism $G\to\operatorname{Sym}(S)$.
\end{remark}

\section{Orbits and transitivity}

\begin{define}[Orbit]\label{def:orbit}
For $s\in S$, its \emph{orbit} is
\[
  \mathcal O_s=G\cdot s=\{g\cdot s:g\in G\}.
\]
An action is \emph{transitive} if $\mathcal O_s=S$ for every
$s\in S$; equivalently, there is only one orbit.
\end{define}

\begin{example}
The natural action $S_n\curvearrowright\{1,\ldots,n\}$ is
transitive: for any $i,j$, a permutation can send $i$ to $j$.
By contrast, $\{1,(1\,2)\}\leq S_3$ has two orbits on
$\{1,2,3\}$, namely $\{1,2\}$ and $\{3\}$.
\end{example}

\begin{example}
For the natural action $\GL_n(\R)\curvearrowright\R^n$,
\[
  \mathcal O_0=\{0\},\qquad
  \mathcal O_{e_1}=\R^n\setminus\{0\}.
\]
An invertible matrix cannot send a nonzero vector to zero.
Conversely, any nonzero $v$ extends to a basis, whose vectors can
be used as the columns of an invertible matrix sending $e_1$ to $v$.
Thus this action has two orbits and is not transitive on $\R^n$.
\end{example}

\begin{proposition}[Orbits partition the set]\label{prop:orbit-partition}
The orbits of an action $G\curvearrowright S$ partition $S$.
In particular, if $S$ is finite, then
\[
  |S|=\sum_{\mathcal O\text{ an orbit}}|\mathcal O|.
\]
\end{proposition}
\begin{proof}
Define $s\sim t$ if $t=g\cdot s$ for some $g\in G$.
This is an equivalence relation: $s=1\cdot s$;
if $t=g\cdot s$, then $s=g^{-1}\cdot t$; and if
$t=g\cdot s$ and $u=h\cdot t$, then $u=(hg)\cdot s$.
The equivalence class of $s$ is exactly $\mathcal O_s$.
\end{proof}

\section{Stabilizers}

\begin{define}[Stabilizer]\label{def:stabilizer}
The \emph{stabilizer} of $s\in S$ is
\[
  G_s=\{g\in G:g\cdot s=s\}.
\]
\end{define}

\begin{lemma}\label{lem:stabilizer-subgroup}
The stabilizer $G_s$ is a subgroup of $G$.
\end{lemma}
\begin{proof}
The identity fixes $s$. If $g,h\in G_s$, then
$(gh)\cdot s=g\cdot(h\cdot s)=s$. Finally, if $g\cdot s=s$,
then $g^{-1}\cdot s=g^{-1}\cdot(g\cdot s)=s$.
\end{proof}

\begin{example}
For $S_n\curvearrowright\{1,\ldots,n\}$,
$(S_n)_n$ consists of the permutations fixing $n$. Restriction to
$\{1,\ldots,n-1\}$ identifies this stabilizer with $S_{n-1}$.
\end{example}

\begin{example}
For $\GL_n(\R)\curvearrowright\R^n$, the stabilizer of $0$ is
all of $\GL_n(\R)$. The stabilizer of $e_1$ consists of the
invertible matrices whose first column is $e_1$, namely those of
block form
\[
  \begin{pmatrix}1&u\\0&B\end{pmatrix},
  \qquad u\in\R^{1\times(n-1)},\quad B\in\GL_{n-1}(\R).
\]
\end{example}

\begin{lemma}[Stabilizers in the same orbit]\label{lem:conjugate-stabilizers}
If $s'=a\cdot s$, then $G_{s'}=aG_sa^{-1}$.
\end{lemma}
\begin{proof}
For $x\in G$,
\begin{align*}
  x\in G_{s'}
  &\iff x\cdot(a\cdot s)=a\cdot s\\
  &\iff (a^{-1}xa)\cdot s=s\\
  &\iff a^{-1}xa\in G_s.
\end{align*}
This is equivalent to $x\in aG_sa^{-1}$.
\end{proof}

\Needspace{29\baselineskip}
\begin{example}[Stabilizer of a short diagonal]\label{ex:short-diagonal-stabilizer}
As in Homework~3, let $D_{12}=\langle x,y\rangle$ act on a regular
hexagon with vertices $v_1,\ldots,v_6$ labeled counterclockwise.
Here $x$ rotates through $\pi/3$ and $y$ reflects across the long
diagonal $\overline{v_1v_4}$. Consider the short diagonal
$d_1=\overline{v_1v_3}$, regarded as an \emph{unordered} pair of
endpoints.

\begin{center}
\begin{tikzpicture}[scale=0.85]
  \foreach \i in {1,...,6} {
    \coordinate (v\i) at ({60*(\i-1)}:1.55);
  }
  \draw[thick,gray] (v1)--(v2)--(v3)--(v4)--(v5)--(v6)--cycle;
  \draw[dashed,thick,Plum] (240:2.05)--(60:2.05);
  \draw[very thick,courseblue] (v1)--(v3)
    node[midway,below left,text=courseblue] {$d_1$};
  \foreach \i in {1,...,6} {\fill (v\i) circle (2pt);}
  \node[right] at (v1) {$v_1$};
  \node[above left] at (v2) {$v_2$};
  \node[above left] at (v3) {$v_3$};
  \node[left] at (v4) {$v_4$};
  \node[below right] at (v5) {$v_5$};
  \node[below right] at (v6) {$v_6$};
  \node[right,text=Plum] at (1.1,1.75) {axis of $x^2y$};
\end{tikzpicture}

{\small The reflection $x^2y$ swaps $v_1$ and $v_3$, preserving $d_1$.}
\end{center}

With subscripts read modulo $6$, rotations and reflections send
\[
  x^k\cdot d_1=d_{1+k},\qquad
  x^ky\cdot d_1=d_{5+k}.
\]
The first equals $d_1$ only for $k=0$, and the second only for
$k=2$. In particular, $x^2y$ preserves $d_1$ by exchanging its
endpoints. Thus
\[
  (D_{12})_{d_1}=\{1,x^2y\},
\]
a subgroup of order $2$. This is the stabilizer that appears in the
short-diagonals part of Homework~3.
\end{example}

\section{The orbit--stabilizer theorem}

\begin{theorem}[Orbit--stabilizer]\label{thm:orbit-stabilizer}
For any action $G\curvearrowright S$ and $s\in S$, there is a
bijection of sets
\[
  G/G_s\longrightarrow\mathcal O_s,
  \qquad gG_s\longmapsto g\cdot s.
\]
Here $G/G_s$ is the set of left cosets; $G_s$ need not be normal.
\end{theorem}
\begin{proof}
If $gG_s=\widetilde gG_s$, then
$\widetilde g^{-1}g\in G_s$ and hence
$g\cdot s=\widetilde g\cdot s$. Thus the map is well defined.
It is surjective by the definition of $\mathcal O_s$.
If $g\cdot s=\widetilde g\cdot s$, then
$\widetilde g^{-1}g\in G_s$, so $gG_s=\widetilde gG_s$.
Therefore it is injective.
\end{proof}

\begin{corollary}\label{cor:orbit-size}
If $G$ is finite, then
\[
  |\mathcal O_s|=[G:G_s],
  \qquad |G|=|\mathcal O_s|\,|G_s|.
\]
In particular, every orbit size divides $|G|$.
\end{corollary}

Together with Proposition~\ref{prop:orbit-partition}, this gives
two useful counting formulas: one measures each orbit using its
stabilizer, and the other adds the sizes of all distinct orbits.
The bijection in Theorem~\ref{thm:orbit-stabilizer} itself does not
require $G$ or $S$ to be finite. For example, the action of
$\GL_n(\R)$ on $\R^n$ has the infinite orbit
$\R^n\setminus\{0\}$ and the singleton orbit $\{0\}$.

\Needspace{16\baselineskip}
\section{Conjugation as an action}

Every group acts on itself by \emph{conjugation}:
\[
  g\cdot x=gxg^{-1}.
\]
Indeed, $1\cdot x=x$ and
$(gh)\cdot x=ghx(gh)^{-1}=g\cdot(h\cdot x)$.
The orbit of $x$ is its \emph{conjugacy class}
$\{gxg^{-1}:g\in G\}$, and its stabilizer is the
\emph{centralizer}
\[
  C_G(x)=\{g\in G:gx=xg\}.
\]
Thus, for finite $G$, orbit--stabilizer gives
\[
  |\operatorname{Cl}_G(x)|=[G:C_G(x)].
\]
Choosing one representative from each conjugacy class gives the
\emph{class equation} $|G|=\sum_x[G:C_G(x)]$, where the sum runs over
those representatives. This action is a starting point for later
counting arguments, including the Sylow theorems.


\end{document}
